\documentclass[a4paper,11pt]{article} % Don't change this
\usepackage{fullpage} % Don't change this

\usepackage{titling} % Don't change this
\pretitle{\vspace{-6em}\begin{center}\larger[2]\bf} % Don't change this
\postauthor{\end{center}} % Don't change this
\preauthor{\vspace{-2em}\begin{center}\smaller[1]} % Don't change this
\postauthor{\end{center}\vspace{-5em}} % Don't change this

\usepackage[shortlabels]{enumitem} % Don't change this
\setlist{itemsep=0em,topsep=0.3em} % Don't change this

%%% Some useful packages.
\usepackage{xcolor}
\usepackage{float}
\usepackage{amsmath,amssymb,mathtools}
\usepackage{graphicx}
\usepackage{subcaption}
\usepackage[hidelinks]{hyperref}
\usepackage{relsize}
\usepackage{cleveref}
\crefname{figure}{Fig.}{Figs.}

\usepackage[backend=bibtex,natbib=true,style=authoryear-icomp,maxnames=3,giveninits=true]{biblatex}
\citetrackerfalse
\bibliography{ref} % put bib entries in ref.bib

\usepackage[linesnumbered]{algorithm2e}
\usepackage{cleveref}
\crefname{equation}{Equation.}{Equations.}
\crefname{figure}{Figure}{Figures.}

\usepackage{lipsum}
\usepackage{wrapfig}
\usepackage{tikz}

%%% Add additional packages or macros below
\usepackage{array}
\usepackage{microtype}
\usepackage[compact]{titlesec}
\titlespacing*{\section}{0pt}{0.8ex plus .2ex}{0.3ex plus .1ex}
\setlength{\abovedisplayskip}{4pt}
\setlength{\belowdisplayskip}{4pt}
\setlength{\abovedisplayshortskip}{2pt}
\setlength{\belowdisplayshortskip}{2pt}
% The title block is three centred environments whose \topsep glue can shrink.
% If page 1 is a hair overfull, TeX shrinks that glue and the title collapses
% onto the author line instead of moving a line to page 2. Remove the shrink.
\raggedbottom
\AtBeginDocument{\setlength{\topsep}{9pt plus 3pt}\setlength{\partopsep}{3pt plus 1pt}}

% Draft aid: red bold TODO markers. Delete this macro's colour (or the macro
% calls) before submission.
\newcommand{\TODO}[1]{\textcolor{red}{\textbf{[TODO: #1]}}}

\newcommand{\Kh}{\widehat{K}}
\newcommand{\Fnorm}[1]{\lVert #1 \rVert_F}
\newcommand{\tgrok}{t_{\mathrm{grok}}}

\begin{document}

\title{Scale or Rotation? Decomposing Neural Tangent Kernel Movement During Grokking}
\author{Group G2 \\
Zain Al-Saffi (s4800836) $\cdot$
Jasper Chong (s4883420) $\cdot$
Aaron D'Mello (s4890257) \\
Zac Kienzle (s4803349) $\cdot$ Samuel Rieger (s4904959)}

\date{}
\maketitle
\TODO{Need to verify all citations by hand}
\section{Topic}

The task is modular addition. The input concatenates the
one-hot codes of $a,b\in\{0,\dots,p-1\}$ and the target is the one-hot code of
$(a+b)\bmod p$. A network $f(x;\theta)$
trained on a fraction of the $p^2$ pairs \emph{groks} when its training loss
falls early but its test loss stays at chance for orders of magnitude longer
before falling abruptly \citep{power2022grokking}. The grokking time $\tgrok$ counts
the training steps from when the training loss first falls below a fixed level
$\tau_{\mathrm{train}}$ to when the test loss first falls below a fixed level
$\tau_{\mathrm{test}}$. The neural tangent kernel (NTK)
of $f$ is
$\Theta_t(x,x')=\langle\nabla_\theta f(x;\theta_t),\nabla_\theta f(x';\theta_t)\rangle$,
and $K_t$ is its empirical version on a fixed probe set of $n$ inputs.
Training is \emph{lazy} when $\Theta_t$ stays at $\Theta_0$, so the network
behaves as its linearisation, and \emph{rich} when $\Theta_t$ moves
\citep{jacot2018ntk,chizat2019lazy}. The leading account of grokking is a
lazy-to-rich transition: a linearised model memorises, but generalising needs
the kernel to move \citep{kumar2024grokking}.

Saying that the kernel moves bundles two claims. A
kernel can change in overall \emph{scale}, with every eigenvalue rescaled and
the eigenbasis fixed, or it can \emph{rotate}, so that the eigenbasis itself
changes. Let $\Kh=K/\Fnorm{K}$ be the normalised kernel, and define the scale
term $S_t=\log(\Fnorm{K_t}/\Fnorm{K_0})$ and the rotation term
$R_t=1-\langle\Kh_t,\Kh_0\rangle_F$. The standard measure of kernel change
mixes the two, since it decomposes as
\begin{equation}
  \Fnorm{K_t-K_0}^2\big/\Fnorm{K_0}^2 = e^{2S_t}+1-2e^{S_t}(1-R_t).
  \label{eq:decomp}
\end{equation}
% : at $R_t=0$ it still reads
% $(1-e^{S_t})^2$, and any exponent fitted to it is unidentified.

The distinction matters because the theorem usually cited to link weight
decay to kernel movement describes pure scaling. A network is $k$-homogeneous
if scaling every parameter by $c$ scales its output by $c^k$; a bias-free ReLU
MLP has $k$ equal to its depth. \citet[Thms.~1--2]{lewkowycz2020l2} show that
an infinitely wide $k$-homogeneous network trained with MSE and $L_2$
strength $\lambda$ has kernel $\Theta_t=e^{-2(k-1)\lambda t}\Theta_0$.
% (as step size $\eta \rightarrow0$, $k\ge2$, conditional on the
% Dyer--Gur-Ari conjecture): eigenvalues decay, eigenvectors static,
In our terms $S_t=-2(k-1)\lambda t$ and $R_t\equiv0$: the kernel shrinks but
never turns, so no features are learned. Yet
\citet[App.~13, arXiv v3]{kumar2024grokking} invoke this theorem to argue that
weight decay ends lazy training, and hence that $\tgrok\propto1/\lambda$, an argument
they call speculative and do not test on the kernel. Rotation is also
necessary: a model that stays in the kernel regime provably needs a constant
fraction of the $p^2$ pairs \citep{mohamadi2024grok}, so at low training
fractions the kernel must turn.

\textbf{Research question.} What sets $\tgrok$? The candidates are kernel
rotation, the scale of the kernel or the weights, and the weight-decay rate.
Three 2026 accounts each back a different one (\S2), and none measures
rotation. We measure all three on one grid and ask which $\tgrok$ tracks.

Our hypothesis H is that rotation is the clock. Let $Y$ be the probe set's
one-hot label matrix and $H=I-\tfrac1n\mathbf{1}\mathbf{1}^\top$ the centring
matrix. The cosine $A_t$ between $HK_tH$ and $HYY^\top H$ is the
centred alignment of \citet{cortes2012alignment}. Rescaling $K_t$ leaves it
unchanged, so only rotation moves it. H says generalisation begins when $A_t$ crosses a
threshold $A_*$, that $A_*$ is shared across $\alpha$, $\eta\lambda$, $N$ and
$p$, and that $\tgrok$ is independent of $S_t$ given $A_t$.

% \textbf{Hypothesis H (rotation clock).} Generalisation begins when task-relevant
% rotation $A_t=\langle HK_tH,\,HYY^\top H\rangle_F$ --- centred alignment
% \citep{cortes2012alignment}, invariant to $S_t$ by construction --- crosses a
% threshold $A_*$; $A_*$ is approximately invariant across $\alpha$, $\eta\lambda$,
% $N$ and $p$, and $\tgrok$ is conditionally independent of $S_t$ given $A_t$.

% \textbf{Discriminating prediction.} In the exact evolution
% $\dot\Theta_t=-2(k-1)\lambda\Theta_t+T_t+T_t^\top$ \citep[Eq.~S5]{lewkowycz2020l2}
% all rotation is carried by $T_t$, which vanishes with width and carries no
% explicit $\lambda$. Fitting $\log\tgrok=c+a\log N+b\log(\eta\lambda)$, our hypothesis predicts
% $a>0,\ b\approx0$ (with $a=1$ if the $O(1/N)$ correction scaling holds); a
% weight-decay clock predicts $a\approx0,\ b\approx-1$. We estimate $a,b$ rather
% than assuming either.

The prediction that separates H from its rivals comes from the exact kernel
evolution $\dot\Theta_t=-2(k-1)\lambda\,\Theta_t+T_t+T_t^\top$
\citep[Eq.~S5, arXiv v1]{lewkowycz2020l2}. The weight-decay term only rescales
$\Theta_t$, so $A_t$ moves only through $T_t$.
As $T_t$ vanishes at infinite width \citep{jacot2018ntk,lewkowycz2020l2},
rotation is a finite-width effect and a wider network turns more slowly. Under
H the target $A_*$ is fixed, so $\tgrok$ is a distance over a rate: it should
grow with $N$ and stay flat in $\eta\lambda$. We fit $\log\tgrok=c+a\log N+b\log(\eta\lambda)$ on a crossed
grid. Under H, $a>0$ and $b\approx0$; under a weight-decay clock, $a\approx0$
and $b\approx-1$; two nonzero exponents mean two clocks, the slower one setting the time.

\section{Significance}

Three 2026 accounts name three different clocks.
\citet{truong2026norm} intervene on the weight norm.
Networks grok when the norm reaches a critical value
$\lVert W\rVert_c\propto p^{0.38}$, and holding the norm at
$\rho\lVert W\rVert_c$ makes $\tgrok$ grow exponentially in $\rho$, at a rate of
about $7.5$ for an MLP and $15$ for unnormalised attention.
\citet{truong2026logit} then shows that under cross-entropy this effect runs
through the logit scale rather than the norm itself, but that this channel is
absent under MSE, where about half of the effect is unexplained.
\citet{kim2026clock} instead makes weight decay the clock, with
$\tgrok\propto(1-\beta)/(\eta\lambda)$. All three are scale-like or
optimiser-level quantities, so none can tell a rotating kernel from a decaying
one; H supplies the missing variable. Our baseline uses MSE and vanilla
gradient descent, where the logit-scale channel is absent, the theorem above
applies as stated, and coupled $L_2$ regularisation coincides with decoupled
weight decay.

The NTK--grokking link, the laziness knob $\alpha$, eNTK spectral tracking,
the norm delay law and norm mediation are all published. What is new here is
decomposition \eqref{eq:decomp} used as a timing variable, the $(a,b)$ test
that separates width-driven rotation from the $\lambda$-clock, and a rotation
account of the norm law under MSE.

\section{Feasibility}

The baseline follows \citet[App.~8.3, arXiv v3]{kumar2024grokking} exactly, so
Tier~0 can be checked against a published figure. It uses one hidden layer of
width $N=100$, MSE loss, full-batch gradient descent at $\eta=100$, $p=23$, and
$90\%$ of the $p^2$ pairs for training. One setting of the swept
hyperparameters is a cell, and the baseline is one cell. A cross-entropy arm
at the baseline cell links our constants to \citet{truong2026norm}, whose
results use cross-entropy.

Each checkpoint records the training and test MSE, the terms $S_t$, $R_t$ and
$A_t$ on the centred kernel, the parameter movement, and $y^\top K^{-1}y$,
which predicts when linearisation breaks. The empirical NTK
on $n=256$ probe points, summed over logits, takes seconds.

Three choices are fixed in advance. First, $\tgrok$ and its thresholds are
pre-registered and sensitivity to $\tau$ is reported. A run has not grokked
if its test loss is above $\tau_{\mathrm{test}}$ when training stops. The run
is kept, and $\tgrok$ is recorded only as longer than the run, which is right
censoring. Second, laziness is set by the centred,
rescaled predictor $\tilde f=\alpha[f(x,\theta)-f(x,\theta_0)]$ with
$\eta=\eta_0/\alpha^2$ \citep{chizat2019lazy}; since $\alpha$, initialisation
scale and label rescaling are one knob, we sweep only $\alpha$. Third,
weight decay is swept in $\eta\lambda$ below the stability bound
$\eta\lambda<2$.

The design has four arms. One is a $6\times6$ grid in
$\alpha\times\eta\lambda$ at the baseline width. One is an $N\times\eta\lambda$
factorial with $N\in\{50,100,400,1600\}$ and
$\eta\lambda\in\{0,10^{-3},10^{-2},10^{-1}\}$ for the $(a,b)$ fit. One is a
clamp arm with $\rho\in\{1,1.25,1.5,1.75,2\}$, and one is a generality arm with
$p\in\{23,29,31\}$ and training fraction in $\{50,70,90\}\%$. Five seeds per
cell, and ten on the clamp arm where the MSE effect is half-size, give about
$350$ runs and roughly $20$ GPU-hours. \TODO{time one run in week~1 and resize;
does $p=23$ grok cleanly? Colab/Kaggle or UQ Bunya?}

The work falls into four tiers, each reportable on its own. Tier~0
reproduces the published $\alpha$ sweep and is a hard gate in week~1, since
later work depends on it. Tier~1 measures $S_t$, $R_t$, $A_t$
and $\tgrok$ over the grid, tests whether $A_*$ is invariant, and exhibits a
cell where kernel movement is large but $A_t$ is flat. Tier~2, the head claim,
is the $(a,b)$ fit with bootstrap intervals and the residuals where a power law
fails. Tier~3, a stretch goal, replicates the clamped-norm law under MSE and
tests whether $A_*$ is invariant to $\rho$.
\TODO{placeholder split} Zac Kienzle builds the infrastructure, Aaron D'Mello
runs Tier~0, Jasper Chong writes the kernel metrics, Samuel Rieger runs the
sweeps and fits, and Zain Al-Saffi does the theory and write-up.

\clearpage
\smaller
\printbibliography

\clearpage
\section*{\textcolor{red}{Draft notes --- delete this page before submission}}
\textcolor{red}{Body is two pages with about one line to spare under
Tectonic 0.17; references are excluded, per the brief, and now start on their
own page. All citations were checked against primary sources on 31~Aug 2026.
Under Tectonic the 31~Aug text ran ten lines onto a third page, so the 4~Sep
rewrite was trimmed to fit. The bib entries were rebuilt in \texttt{ref.bib} and
rechecked against the arXiv abstract pages on 4~Sep 2026. In a later 4~Sep
pass the run-in bold headings were removed from every paragraph except the
research question, the research question was rewritten in full sentences, and
a few phrases were shortened so the body still fits. On 6~Sep the grokking
time was spelled out as a count of steps between two threshold crossings, and
``cell'', ``has not grokked'' and ``right censoring'' were each defined in
plain words at first use. To pay for those lines the research question was
tightened (``the grokking time'' dropped after $\tgrok$, ``itself'' dropped,
``rotation of the kernel'' became ``kernel rotation''), the two compute TODOs
were merged into one, ``gradient-descent stability bound'' became ``stability
bound'', and ``not dropped'' was folded into ``the run is kept''. Page~2 is
now full with no line to spare.}
\begin{enumerate}[1.]
  \item \textcolor{red}{\textbf{What changed in this revision.} The hypothesis is now a \emph{threshold} claim ($A_*$ invariant, $\tgrok \perp S_t \mid A_t$) plus a \emph{rate-law} claim tested by estimating $(a,b)$ in $\log\tgrok = c + a\log N + b\log(\eta\lambda)$. That regression, not the old mediation test, is the head claim: it is the one experiment that separates the rotation clock, the scale clock and \citeauthor{kim2026clock}'s weight-decay clock, and it fails informatively (both exponents nonzero $\Rightarrow$ two clocks). Tier~3 was demoted to a stretch replication because \citet{truong2026logit} already ran a norm mediation test under cross-entropy.}
  \item \textcolor{red}{\textbf{Verified numbers} in \citet{truong2026norm}: $\lVert W\rVert_c=\{42,48,55,61,67\}$ for $p=\{29,41,59,79,97\}$, $\lVert W\rVert_c\propto p^{0.38}$; $\tgrok\propto e^{\alpha\rho}$ with $\alpha\approx7.4$--$7.5$ (MLP) against ${\approx}15$ (unnormalised attention), $R^2>0.99$; ``the form of the delay law is cross-architectural while its rate is not''; LayerNorm removes the dependence. Their loss is cross-entropy --- hence our CE arm. \citet{truong2026logit}: logit-scale mediation $R^2=0.97$ under CE, channel absent under MSE, norm effect ${\approx}$ half size.}
  \item \textcolor{red}{\textbf{Verified as stated:} Lewkowycz \& Gur-Ari Thm.~1 and Thm.~2 (eigenvalues decay, eigenvectors static; assumes MSE, $k\ge2$, gradient flow, infinite width, Dyer--Gur-Ari conjecture), Eq.~(S5) with $T_t$ from (S6) in arXiv v1 (these are (S8) and (S9) in v2, so the body cites the version); Kumar App.~8.3 config and App.~13 for weight decay (an earlier note said 14, but arXiv v3 and the ICLR camera-ready both number it 13, and App.~14 is momentum); Mohamadi et al.\ (ICML 2024, PMLR~235); Atanasov et al.; Gromov's quadratic-activation closed form; Cortes et al., JMLR~13, 795--828.  \textbf{Venue lines rechecked 4~Sep 2026:} Jacot (NeurIPS~2018), Chizat (NeurIPS~2019) and Lewkowycz (NeurIPS~2020) appear in the NeurIPS proceedings index; Nanda appears on OpenReview as ICLR~2023 (forum 9XFSbDPmdW); Power appears on the ICLR 2021 MATH-AI workshop page, while the arXiv posting is dated 2022. Kumar's ICLR~2024 venue rests on the camera-ready checked on 31~Aug.}
  \item \textcolor{red}{\textbf{Open decisions.} Coordinator sign-off (scope email to Dr Nan Ye before 14 Sep); AI-use statement here or in the report; whether ``pre-registered'' means public (OSF) or fixed in the repo before sweeps run; whether the CE arm is one cell or a full $\alpha\times\eta\lambda$ replication; whether to keep the $p$ and training-fraction arms at all, given they buy generality rather than discrimination.}
  \item \textcolor{red}{\textbf{Held for the 12-page report:} the derivation of \eqref{eq:decomp}; the tier table; uncentred-alignment and first-logit eNTK robustness variants; the censoring model; $y^\top K^{-1}y=\lVert w-w_0\rVert^2$ near the kernel limit; the confound demonstration figure. An abstract would cost about six lines here.}
  \item \textcolor{red}{\textbf{Dropped from the body for space on 4~Sep 2026} (restore if lines are freed): ``and cannot tell them apart''; ``and so leaves $A_t$ unchanged; alignment moves only through $T_t$'' (now ``so $A_t$ moves only through $T_t$''); ``and the usual summary statistic does not separate the two''; ``each observes only a magnitude, which rescaling moves but rotation does not'' (the point is made in \S2); ``while LayerNorm removes the dependence entirely''; ``and the $(a,b)$ fit separates the three accounts''; ``Since $A_t$ is computable online and is not flat during memorisation, it is also a progress measure in the sense of \citet{nanda2023progress}''; ``The kernel evolution equation above is exact, so Tiers~1--2 stand regardless''; ``Runs take minutes on one GPU''; ``with $a=1$ if the finite-width correction is $O(1/N)$''; ``so the fits are not selection-biased''; ``where the published mediator does not exist''; ``The dataset is $(a+b)\bmod p$''; ``or $\lambda<0.02$ at $\eta=100$''; ``computed with \texttt{torch.func}''; the pilot-figure placeholder and its caption ``Baseline grokking with $S_t$ and $A_t$ overlaid'' (the figure was dropped on 4~Sep at the group's request). \textbf{Tier~4 moved here from the work plan:} droppable theory that bounds $\lVert T_t\rVert$ for the two-layer quadratic-activation network ($k=3$), whose grokked solution is closed-form \citep{gromov2023grokking}, giving an analytic $A_\infty$ and $\tgrok\sim A_\infty/\dot A$; the kernel evolution equation (S5/S8) is exact, so Tiers~1--2 stand regardless. In the second 4~Sep pass (flow and set-up) the pilot figure was removed at the group's request, ``Why this regime'' was merged into the end of the dispute paragraph and its clause ``so the distinction \citeauthor{kim2026clock} draws between coupled and decoupled weight decay does not arise'' was dropped, and the Novelty list lost ``the MLP-versus-attention comparison''. Restoring Tier~4 costs about three lines; \texttt{\textbackslash linespread\{0.97\}} would free roughly that many if the group prefers it to cutting text.}
  \item \textcolor{red}{\textbf{Template.} Font size, page size and the ``don't change'' lines untouched; added \texttt{array}, \texttt{microtype} (character protrusion only under XeTeX), \texttt{titlesec} (tighter section spacing), tighter display skips, \texttt{\textbackslash raggedbottom} and shrink-free \texttt{\textbackslash topsep} (without this, a page-1 overrun of a few points shrinks the title block onto the author line instead of moving a line). Reverting lines~39--45 overflows by several lines. Submit as \emph{G2\_proposal\_\textless SubmitterName\textgreater.pdf}.}
\end{enumerate}

\end{document}
